% GENERATED FILE. Edit canonical lessons, then run npm run generate:books.
% canonical-source-hash: fnv1a64:7723a2d9bc008d16
% canonical-lessons: BN-C388-tram, BN-C388-metro, BN-C388-oto, BN-C388-bhen, BN-C388-dourano

\chapter{Trams, the Metro and Autos}
\label{ch:bn-trams-the-metro-and-autos}
\hlchaptermodality{388}

\begin{chapteropening}
\textbf{By the end of this chapter:} \emph{I can choose between a tram, the metro and an auto, and run for the bus.}

Complete the last lesson of chapter 388: I can choose between a tram, the metro and an auto, and run for the bus.
\end{chapteropening}

\section[\texorpdfstring{\bn{ট্রাম} (ṭrām)}{ট্রাম (ṭrām)}]{\bn{ট্রাম} (ṭrām) — a tram}

\label{lesson:BN-C388-tram}



\begin{quote}
Before the new one: say the Bengali for to commute, then the Bengali for return (of a journey).
\end{quote}

\subsection*{You'll want to know: \bn{ট্রাম}}
\textbf{\bn{ট্রাম}} — \emph{ṭrām} — \textquotedblleft{}a tram\textquotedblright{}.

\begin{quote}
\bn{ট্রামটা} \bn{ধীরে} \bn{চলে}, \bn{কিন্তু} \bn{ভাড়া} \bn{সস্তা।}
\end{quote}

\begin{quote}
\emph{ṭrāmṭā dhīre chôle, kintu bhāṛā sôstā.}
\end{quote}

\begin{quote}
The tram is slow, but the fare is cheap.
\end{quote}

\begin{cousinweb}
\textbf{\bn{ট্রাম}} (\emph{ṭrām}) is borrowed from English \emph{tram}, short for \emph{tramway}.
\end{cousinweb}

\begin{grammarlens}[title={What you've built}]
One more word for getting around.
\end{grammarlens}

\subsection*{Your turn}
\begin{itemize}
\raggedright
  \item \emph{Say it:} \emph{ṭrām}
  \item \emph{Say it:} \emph{ṭrām}, once more
  \item \emph{Say it:} \emph{ṭrāmṭā dhīre chôle, kintu bhāṛā sôstā.} — The tram is slow, but the fare is cheap.
  \item \emph{Recall:} say the Bengali for to commute, then the Bengali for return (of a journey), then say \emph{ṭrām} again
\end{itemize}

\begin{tcolorbox}[breakable,colback=teal!4,colframe=teal!35!black,title={Before you move on}]
What does \bn{ট্রাম} mean? (\textquotedblleft{}A tram\textquotedblright{}.) Say it once more.
\end{tcolorbox}
\section[\texorpdfstring{\bn{মেট্রো} (meṭro)}{মেট্রো (meṭro)}]{\bn{মেট্রো} (meṭro) — the metro; an underground railway}

\label{lesson:BN-C388-metro}



\begin{quote}
Before the new one: say the Bengali for return (of a journey), then the Bengali for a tram.
\end{quote}

\subsection*{You'll want to know: \bn{মেট্রো}}
\textbf{\bn{মেট্রো}} — \emph{meṭro} — \textquotedblleft{}the metro; an underground railway\textquotedblright{}.

\begin{quote}
\bn{মেট্রো} \bn{ট্রামের} \bn{চেয়ে} \bn{দ্রুত।}
\end{quote}

\begin{quote}
\emph{meṭro ṭrāmer cheye drutô.}
\end{quote}

\begin{quote}
The metro is faster than the tram.
\end{quote}

\bn{মেট্রোতে}, by metro, on the metro.

\begin{cousinweb}
\textbf{\bn{মেট্রো}} (\emph{meṭro}) is borrowed from English \emph{metro}, short for \emph{metropolitan railway}.
\end{cousinweb}

\begin{grammarlens}[title={What you've built}]
One more word for getting around.
\end{grammarlens}

\subsection*{Your turn}
\begin{itemize}
\raggedright
  \item \emph{Say it:} \emph{meṭro}
  \item \emph{Say it:} \emph{meṭro}, once more
  \item \emph{Say it:} \emph{meṭro ṭrāmer cheye drutô.} — The metro is faster than the tram.
  \item \emph{Recall:} say the Bengali for return (of a journey), then the Bengali for a tram, then say \emph{meṭro} again
\end{itemize}

\begin{tcolorbox}[breakable,colback=teal!4,colframe=teal!35!black,title={Before you move on}]
What does \bn{মেট্রো} mean? (\textquotedblleft{}The metro; an underground railway\textquotedblright{}.) Say it once more.
\end{tcolorbox}
\section[\texorpdfstring{\bn{অটো} (ôṭo)}{অটো (ôṭo)}]{\bn{অটো} (ôṭo) — an autorickshaw}

\label{lesson:BN-C388-oto}



\begin{quote}
Before the new one: say the Bengali for a tram, then the Bengali for the metro.
\end{quote}

\subsection*{You'll want to know: \bn{অটো}}
\textbf{\bn{অটো}} — \emph{ôṭo} — \textquotedblleft{}an autorickshaw\textquotedblright{}.

\begin{quote}
\bn{স্টেশন} \bn{থেকে} \bn{আমরা} \bn{অটোতে} \bn{বাড়ি} \bn{ফিরলাম।}
\end{quote}

\begin{quote}
\emph{sṭeshôn theke āmrā ôṭote bāṛi phirlām.}
\end{quote}

\begin{quote}
We came home from the station by auto.
\end{quote}

The three-wheeler with a motor; a \bn{রিকশা} is pulled or pedalled.

\begin{cousinweb}
\textbf{\bn{অটো}} (\emph{ôṭo}) is borrowed from English \emph{auto}, clipped from \emph{autorickshaw}; the first letter is \bn{অ}, said as an open ô.
\end{cousinweb}

\begin{grammarlens}[title={What you've built}]
One more word for getting around.
\end{grammarlens}

\subsection*{Your turn}
\begin{itemize}
\raggedright
  \item \emph{Say it:} \emph{ôṭo}
  \item \emph{Say it:} \emph{ôṭo}, once more
  \item \emph{Say it:} \emph{sṭeshôn theke āmrā ôṭote bāṛi phirlām.} — We came home from the station by auto.
  \item \emph{Recall:} say the Bengali for a tram, then the Bengali for the metro, then say \emph{ôṭo} again
\end{itemize}

\begin{tcolorbox}[breakable,colback=teal!4,colframe=teal!35!black,title={Before you move on}]
What does \bn{অটো} mean? (\textquotedblleft{}An autorickshaw\textquotedblright{}.) Say it once more.
\end{tcolorbox}
\section[\texorpdfstring{\bn{ভ্যান} (bhæn)}{ভ্যান (bhæn)}]{\bn{ভ্যান} (bhæn) — a van; a cycle van}

\label{lesson:BN-C388-bhen}



\begin{quote}
Before the new one: say the Bengali for the metro, then the Bengali for an autorickshaw.
\end{quote}

\subsection*{You'll want to know: \bn{ভ্যান}}
\textbf{\bn{ভ্যান}} — \emph{bhæn} — \textquotedblleft{}a van; a cycle van\textquotedblright{}.

\begin{quote}
\bn{একটা} \bn{ভ্যান} \bn{আমাদের} \bn{মালপত্র} \bn{স্টেশনে} \bn{নিয়ে} \bn{গেল।}
\end{quote}

\begin{quote}
\emph{ekṭā bhæn āmāder mālpôtrô sṭeshône niye gelô.}
\end{quote}

\begin{quote}
A van took our luggage to the station.
\end{quote}

In a small town a \bn{ভ্যান} is often a pedalled tricycle with a flat wooden bed.

\begin{cousinweb}
\textbf{\bn{ভ্যান}} (\emph{bhæn}) is borrowed from English \emph{van}; Bengali has no v, so the word begins with \bn{ভ}, and the y-phala gives the English vowel æ.
\end{cousinweb}

\begin{grammarlens}[title={What you've built}]
One more word for getting around.
\end{grammarlens}

\subsection*{Your turn}
\begin{itemize}
\raggedright
  \item \emph{Say it:} \emph{bhæn}
  \item \emph{Say it:} \emph{bhæn}, once more
  \item \emph{Say it:} \emph{ekṭā bhæn āmāder mālpôtrô sṭeshône niye gelô.} — A van took our luggage to the station.
  \item \emph{Recall:} say the Bengali for the metro, then the Bengali for an autorickshaw, then say \emph{bhæn} again
\end{itemize}

\begin{tcolorbox}[breakable,colback=teal!4,colframe=teal!35!black,title={Before you move on}]
What does \bn{ভ্যান} mean? (\textquotedblleft{}A van; a cycle van\textquotedblright{}.) Say it once more.
\end{tcolorbox}
\section[\texorpdfstring{\bn{দৌড়ানো} (douṛāno)}{দৌড়ানো (douṛāno)}]{\bn{দৌড়ানো} (douṛāno) — to race; to run fast}

\label{lesson:BN-C388-dourano}



\begin{quote}
Before the new one: say the Bengali for an autorickshaw, then the Bengali for a van.
\end{quote}

\subsection*{You'll want to know: \bn{দৌড়ানো}}
\textbf{\bn{দৌড়ানো}} — \emph{douṛāno} — \textquotedblleft{}to race; to run fast\textquotedblright{}.

\begin{quote}
\bn{অরুণ} \bn{দৌড়ে} \bn{বাসে} \bn{উঠল।}
\end{quote}

\begin{quote}
\emph{ôrun douṛe bāse uṭhlô.}
\end{quote}

\begin{quote}
Arun ran and got on the bus.
\end{quote}

Past (he/she …): \bn{তিনি} \bn{দৌড়ালেন}, \emph{tini douṛālen}.

The -\bn{এ} form \bn{দৌড়ে} is the one you hear most: \bn{দৌড়ে} \bn{এল}, came running.

The first vowel sign is \bn{ৌ}, the vowel ou.

\begin{cousinweb}
\textbf{\bn{দৌড়ানো}} (\emph{douṛāno}) is inherited: CDIAL lists Bengali \emph{dauṛā} and \emph{dauṛāna}, \textquotedblleft{}to run\textquotedblright{}, under Sanskrit \emph{drávati}, \textquotedblleft{}runs\textquotedblright{} (CDIAL 6624), and notes that the word probably came east from the languages to the west.
\end{cousinweb}

\begin{grammarlens}[title={What you've built}]
One more word for getting around.
\end{grammarlens}

\subsection*{Your turn}
\begin{itemize}
\raggedright
  \item \emph{Say it:} \emph{douṛāno}
  \item \emph{Say it:} \emph{douṛāno}, once more
  \item \emph{Say it:} \emph{ôrun douṛe bāse uṭhlô.} — Arun ran and got on the bus.
  \item \emph{Recall:} say the Bengali for an autorickshaw, then the Bengali for a van, then say \emph{douṛāno} again
\end{itemize}

\begin{tcolorbox}[breakable,colback=teal!4,colframe=teal!35!black,title={Before you move on}]
What does \bn{দৌড়ানো} mean? (\textquotedblleft{}To race; to run fast\textquotedblright{}.) Say it once more.
\end{tcolorbox}
